\documentclass[11pt]{article}
\usepackage[margin=1in]{geometry}
\usepackage{amsmath,amssymb}
\usepackage{xcolor}
\usepackage{booktabs}
\usepackage{enumitem}
\usepackage{fancyhdr}
\usepackage{tcolorbox}
\tcbuselibrary{skins,breakable}
\usepackage{hyperref}
\hypersetup{colorlinks=true, linkcolor=blue, urlcolor=blue}

\definecolor{primary}{RGB}{30,64,140}
\definecolor{tipbg}{RGB}{235,244,255}
\definecolor{examplebg}{RGB}{240,248,235}
\definecolor{formulabg}{RGB}{255,247,224}
\definecolor{warnbg}{RGB}{255,236,236}

\pagestyle{fancy}
\fancyhf{}
\renewcommand{\headrulewidth}{0.4pt}
\cfoot{\thepage}

\newtcolorbox{tipbox}[1][Tip]{
  colback=tipbg, colframe=primary, fonttitle=\bfseries,
  title=#1, breakable, rounded corners
}
\newtcolorbox{examplebox}[1][Example]{
  colback=examplebg, colframe=primary!70!black, fonttitle=\bfseries,
  title=#1, breakable, rounded corners
}
\newtcolorbox{formulabox}[1][Formula]{
  colback=formulabg, colframe=primary!70!black, fonttitle=\bfseries,
  title=#1, breakable, rounded corners
}
\newtcolorbox{warnbox}[1][Warning]{
  colback=warnbg, colframe=red!60!black, fonttitle=\bfseries,
  title=#1, breakable, rounded corners
}


\title{\color{primary}\textbf{Calendar --- Practice Questions}\\[4pt]\large Medium to Difficult, With Full Solutions}
\author{Aptitude Grind}
\date{}

\lhead{\small Calendar --- Practice Questions}
\rhead{\small Aptitude Grind}

\begin{document}
\maketitle
\thispagestyle{fancy}

\noindent\textit{Difficulty is marked as \textbf{[M]} Medium or \textbf{[H]} Hard. Solve on your own
first, then check the solution. Solutions use the methods from \texttt{01-fundamentals.tex},
\texttt{02-fast-tricks-and-shortcuts.tex}, and \texttt{03-question-pattern-catalog.tex}. Every
numeric answer below was cross-checked independently against Python's \texttt{datetime}/\texttt{calendar}
standard library before being finalized.}

\bigskip
\hrule
\bigskip

\subsection*{Q1 [M] --- Direct Date Lookup (Pattern A1)}
What was the day of the week on 26 January 1950?

\begin{tipbox}[Solution]
Reference: 1 January 1900 was a \textbf{Monday}.
\begin{itemize}[leftmargin=1.4em]
  \item Years 1900--1949 (up to start of 1950): count odd days for years 1900 to 1949 inclusive.
  \item Leap years 1900--1949: 1904, 1908, \ldots, 1948 --- but 1900 is \emph{not} leap (century rule)
    $\Rightarrow$ 12 leap years, 38 ordinary years.
  \item Odd days $= (12 \times 2) + (38 \times 1) = 24 + 38 = 62 \equiv 62 - 56 = 6 \pmod 7$
  \item Monday $+ 6$ odd days $\to$ \textbf{Sunday} (this gives the day for 1 Jan 1950)
  \item Add 25 more days to reach 26 Jan: $25 \bmod 7 = 4$ odd days
  \item Sunday $+ 4 \to$ \textbf{Thursday}
\end{itemize}
\textbf{Answer: Thursday}
\end{tipbox}

\subsection*{Q2 [M] --- N Days After (Pattern B1)}
Today is Wednesday. What day will it be after 90 days?

\begin{tipbox}[Solution]
$90 \bmod 7 = 90 - 84 = 6$ odd days. \\
Wednesday $+ 6$ days $\to$ \textbf{Tuesday}.\\
\textbf{Answer: Tuesday}
\end{tipbox}

\subsection*{Q3 [M] --- Nth Weekday of a Month (Pattern C1)}
What date is the fourth Thursday of November 2026?

\begin{tipbox}[Solution]
1 November 2026 is a \textbf{Monday}. The first Thursday is therefore the 4th (Mon$\to$Thu is 3
days). Each subsequent Thursday is 7 days later:
\[
\text{1st} = 4,\quad \text{2nd} = 11,\quad \text{3rd} = 18,\quad \text{4th} = 25
\]
\textbf{Answer: 25 November 2026}
\end{tipbox}

\begin{warnbox}[Why This Is a Common Trap]
People often assume the Nth weekday is always $7\times(N-1)+1$ (i.e. assume the 1st of the month IS
the target weekday). Always find the \emph{first occurrence} of the target weekday first --- it
usually isn't the 1st --- then add multiples of 7.
\end{warnbox}

\subsection*{Q4 [M] --- Last Weekday of a Month (Pattern C2)}
What is the date of the last Sunday of September 2026? (September has 30 days.)

\begin{tipbox}[Solution]
1 September 2026 is a \textbf{Tuesday}, so the Sundays in September 2026 fall on 6, 13, 20, 27. The
last one before the month ends on the 30th is the \textbf{27th}.\\
\textbf{Answer: 27 September 2026}
\end{tipbox}

\subsection*{Q5 [M] --- Days Between Two Dates (Pattern F1)}
How many days are there between 15 August 2010 and 25 December 2010 (counting from the first date to
the second, exclusive of the start date)?

\begin{tipbox}[Solution]
Days remaining in August after the 15th: $31-15=16$ \\
Add full months: September (30) + October (31) + November (30) = 91 \\
Add days into December: 25 \\
Total: $16+91+25 = \textbf{132 days}$
\end{tipbox}

\subsection*{Q6 [M] --- Same Date, Different Year (Pattern F3)}
If 15 March 2015 was a Sunday, what day was 15 March 2016?

\begin{tipbox}[Solution]
2015 is an ordinary year, but the gap 15 Mar 2015 $\to$ 15 Mar 2016 \textbf{includes 29 Feb 2016}
(since 2016 is a leap year and the leap day falls before 15 March). So this interval has
\textbf{2 odd days}, not 1. Sunday $+2 \to$ \textbf{Tuesday}.\\
\textbf{Answer: Tuesday}
\end{tipbox}

\begin{warnbox}[Why This Trips People Up]
The ``add 1 odd day per year'' shortcut only works when the year in question is ordinary
\textbf{and} the leap day (if any) doesn't fall between the two matching dates. Always check whether
29 February lies inside the specific date range you're bridging --- not just whether the calendar
year is leap.
\end{warnbox}

\subsection*{Q7 [H] --- Weekdays Occurring Five Times (Pattern C3)}
March 2026 has 31 days and 1 March 2026 is a Sunday. Which weekdays occur five times that month?

\begin{tipbox}[Solution]
$L = 31-28=3$, so the starting weekday and the next two occur five times. Starting weekday is Sunday
$\Rightarrow$ \textbf{Sunday, Monday, Tuesday}.\\
\textbf{Answer: Sunday, Monday, Tuesday}
\end{tipbox}

\subsection*{Q8 [H] --- Counting Leap Years in a Range (Pattern E2)}
How many leap years are there from 1601 to 2000 (inclusive)?

\begin{tipbox}[Solution]
Using $\lfloor N/4\rfloor - \lfloor N/100\rfloor + \lfloor N/400\rfloor$:
\begin{align*}
\text{count}(2000) &= 500-20+5=485\\
\text{count}(1600) &= 400-16+4=388\\
\text{Leap years in }[1601,2000] &= 485-388 = \textbf{97}
\end{align*}
\end{tipbox}

\subsection*{Q9 [H] --- Reverse Problem: Nth Occurrence to Other Occurrences (Pattern G1)}
The third Saturday of a certain month falls on the 16th. On what date does the first Saturday fall,
and on what date does the last Saturday fall if the month has 31 days?

\begin{tipbox}[Solution]
\textbf{First Saturday:} subtract 2 weeks from the 3rd occurrence: $16-14=\textbf{2nd}$.\\
\textbf{Remaining Saturdays:} 2, 9, 16, 23, and (since the month has 31 days) 30 also fits
($23+7=30\le31$), so the \textbf{last Saturday is the 30th}.\\
\textbf{Answer:} First Saturday = 2nd; last Saturday = 30th
\end{tipbox}

\subsection*{Q10 [H] --- Same Calendar Year (Pattern D1)}
The calendar for the year 2023 will be the same as the calendar for which future year?

\begin{tipbox}[Solution]
2023 is ordinary. Track running odd days year by year starting the count from 2023$\to$2024:
\begin{center}
\begin{tabular}{lcc}
\toprule
\textbf{Transition} & \textbf{Odd days added} & \textbf{Running total (mod 7)} \\
\midrule
2023 $\to$ 2024 (ordinary) & 1 & 1 \\
2024 $\to$ 2025 (2024 leap) & 2 & 3 \\
2025 $\to$ 2026 (ordinary) & 1 & 4 \\
2026 $\to$ 2027 (ordinary) & 1 & 5 \\
2027 $\to$ 2028 (ordinary) & 1 & 6 \\
2028 $\to$ 2029 (2028 leap) & 2 & 1 \\
2029 $\to$ 2030 (ordinary) & 1 & 2 \\
2030 $\to$ 2031 (ordinary) & 1 & 3 \\
2031 $\to$ 2032 (ordinary) & 1 & 4 \\
2032 $\to$ 2033 (2032 leap) & 2 & 6 \\
2033 $\to$ 2034 (ordinary) & 1 & 0 \\
\bottomrule
\end{tabular}
\end{center}
Running total first hits a multiple of 7 at \textbf{2034}, and 2034 is also an ordinary year
(matching 2023's type).\\
\textbf{Answer: 2034}
\end{tipbox}

\begin{warnbox}[Why This Is a Common Trap]
Coaching material often claims ordinary years always repeat after ``6, 11, 17, or 28'' years. The 17
figure doesn't actually occur for any real year when checked directly --- don't assume a fixed gap;
always run the odd-day total.
\end{warnbox}

\subsection*{Q11 [H] --- Century Reasoning}
Prove that the last day of a century (year ending in 00) can never be a Tuesday, Thursday, or
Saturday.

\begin{tipbox}[Solution]
The last day of the n-th century corresponds to a cumulative odd-day count equal to $5n \bmod 7$ for
n not a multiple of 4, and $0$ when n is a multiple of 4 (from the 400-year rule). Evaluate
$5n \bmod 7$ for n = 1, 2, 3 (since n=4 gives 0):
\[
n{=}1: 5 \qquad n{=}2: 10\equiv 3 \qquad n{=}3: 15\equiv 1 \qquad n{=}4: 0
\]
So the only possible odd-day residues for a century's last day are $\{0,1,3,5\}$, which map to only
four of the seven weekdays. Direct computation confirms centuries can only end on
\textbf{Sunday, Monday, Wednesday, or Friday} --- never Tuesday, Thursday, or Saturday.\\
\textbf{Answer:} Because only 4 distinct odd-day residues (0, 1, 3, 5) are possible for century
boundaries, only 4 of the 7 weekdays are reachable --- and Tuesday, Thursday, Saturday are never
among them.
\end{tipbox}

\subsection*{Q12 [H] --- Reverse Problem (Multi-Clue Logic)}
1 January of a certain year is a Friday. 1 January of the following year is a Sunday. What can you
conclude about the certain year, and what day is 1 January two years later?

\begin{tipbox}[Solution]
Going from Friday to Sunday is a shift of \textbf{2 odd days} $\Rightarrow$ the certain year must be
a \textbf{leap year} (only leap years contribute 2 odd days). So the certain year has 366 days, and
the following year is ordinary, contributing 1 odd day. Sunday $+1 \to$ \textbf{Monday} for 1
January two years later.\\
\textbf{Answer:} The certain year is a leap year; 1 January two years later is a Monday.
\end{tipbox}

\subsection*{Q13 [H] --- Working-Day Counting (Pattern H2)}
An office is closed every Sunday. A 31-day month starts on a Saturday. How many working days does
the month have?

\begin{tipbox}[Solution]
$L = 31-28=3$, so the starting weekday and the next two occur five times: \textbf{Saturday, Sunday,
Monday}. Sunday is among them $\Rightarrow$ Sunday occurs \textbf{5} times in this month. \\
Working days $=$ total days $-$ Sundays $= 31-5 = \textbf{26}$.\\
\textbf{Answer: 26 working days}
\end{tipbox}

\begin{warnbox}[Why This Is a Common Trap]
Assuming every month has exactly 4 Sundays. Whether Sunday occurs 4 or 5 times depends entirely on
the month's length and starting weekday --- check both before subtracting.
\end{warnbox}

\subsection*{Q14 [H] --- Century-Spanning Problem}
If 1 January 2001 was a Monday, what day of the week is 1 January 2101?

\begin{tipbox}[Solution]
We need the total odd days contributed by the years 2001 through 2100 (100 full years elapsing).

\textbf{Leap years in 2001--2100:} every year divisible by 4 is a leap year \emph{unless} it's a
century year not divisible by 400. In this range: 2004, 2008, \ldots, 2096 are leap
$\Rightarrow$ $(2096-2004)/4+1=24$ leap years. The only century year in range, 2100, is \textbf{not}
divisible by 400, so it is \textbf{not} leap --- it counts as an \textbf{ordinary} year instead.
\begin{itemize}[leftmargin=1.4em]
  \item Leap years: 24
  \item Ordinary years: $100-24=76$
  \item Odd days $= (24\times2)+(76\times1)=48+76=124$
  \item $124 \bmod 7 = 124-119=5$
\end{itemize}
Monday $+5$ odd days $\to$ \textbf{Saturday}.\\
\textbf{Answer: 1 January 2101 is a Saturday.}
\end{tipbox}

\begin{warnbox}[Why This Is a Common CAT/Placement Trap]
The trap is treating 2100 as a normal leap year because it's divisible by 4. Since it's a century
year \textbf{not} divisible by 400, it must be excluded --- this single oversight shifts the final
answer by one full odd day (which changes the answer from Saturday to Sunday if missed).
\end{warnbox}

\section*{Answer Key Summary}
\begin{center}
\begin{tabular}{ll}
\toprule
\textbf{Question} & \textbf{Answer} \\
\midrule
Q1 & Thursday \\
Q2 & Tuesday \\
Q3 & 25 November 2026 \\
Q4 & 27 September 2026 \\
Q5 & 132 days \\
Q6 & Tuesday \\
Q7 & Sunday, Monday, Tuesday \\
Q8 & 97 \\
Q9 & First = 2nd; Last = 30th \\
Q10 & 2034 \\
Q11 & Only Sun/Mon/Wed/Fri possible for century-end (proof-based) \\
Q12 & Certain year is leap; 2 years later = Monday \\
Q13 & 26 working days \\
Q14 & Saturday \\
\bottomrule
\end{tabular}
\end{center}

\end{document}
